\clearpage
\section{Implementation}

This chapter maps the design decisions to the \textbf{files that are actually
delivered}. The goal is to show where the clustered behavior is implemented
and what was added with respect to \textbf{assignment 3}.

\subsection{Implementation map}

\begin{center}
\begin{tabular}{|>{\bfseries}p{3.2cm}|p{10.2cm}|}
    \hline
    Package or class & Implementation responsibility \\
    \hline
    \texttt{pcd.shas.Main}
        & Starts one interactive cluster node and selects the local actor startup based on the CLI role. \\
    \hline
    \texttt{pcd.shas.DemoMain}
        & Starts three separate JVM processes and drives a visible distributed scenario for manual demonstration. \\
    \hline
    \texttt{pcd.shas.runtime.NodeStartup}
        & Parses startup arguments and builds the Pekko Cluster configuration for host, port, seed nodes, and cluster role. \\
    \hline
    \texttt{pcd.shas.controlunit.ControlUnitActor}
        & Implements the alarm state machine, the recovery logic, receptionist integration, and cluster-singleton initialization. \\
    \hline
    \texttt{pcd.shas.keypad.KeypadActor}
        & Collects user input, discovers the control unit, and forwards PIN or arming requests. \\
    \hline
    \texttt{pcd.shas.sensor.SensorActor}
        & Represents one distributed sensor with a stable \texttt{sensorId} and forwards activations as typed messages. \\
    \hline
    \texttt{pcd.shas.siren.SirenActor}
        & Implements the siren state, registers itself for discovery, and reacts to activation or deactivation commands. \\
    \hline
    \texttt{pcd.shas.common}
        & Contains shared serializable messages and enums exchanged across cluster nodes. \\
    \hline
\end{tabular}
\end{center}

\subsection{What changed from assignment 3}

The main technical differences are concentrated in a small number of places:
\begin{itemize}
    \item \textbf{local startup} became \textbf{role-based cluster startup};
    \item the local control unit became a \textbf{cluster singleton};
    \item direct local references were replaced by \textbf{receptionist discovery};
    \item the state machine gained the extra state \texttt{STARTUP\_RECOVERY};
    \item exchanged messages were made explicitly \textbf{serializable}.
\end{itemize}

The rest of the alarm logic was intentionally kept close to assignment 3, so
that the distributed mechanisms stay easy to identify and explain.

\subsection{Node startup and cluster roles}

The entry point is \textbf{role-based}. \texttt{Main} accepts \texttt{control-unit},
\texttt{keypad}, and \texttt{sensor} as startup roles and launches exactly one
interactive node per JVM. \texttt{NodeStartup} converts CLI arguments into a
small node description and then builds the final Pekko configuration overlay.

This keeps the \textbf{distributed deployment explicit}: each JVM joins the same
cluster, but each one hosts only the actors needed for its assigned role.

\subsection{Control unit as cluster singleton}

The alarm state must have a \textbf{single owner}. For this reason, the control unit is
initialized as a Pekko \emph{cluster singleton}. Multiple nodes may have the
\texttt{control-unit} cluster role, but only one active \texttt{ControlUnitActor}
exists at a time.

This is the most important architectural change with respect to assignment 3.
The non-clustered version could rely on one local control actor by
construction; the clustered version must enforce that property explicitly.
If the hosting node fails, Pekko can recreate the singleton on another
eligible node, and the recreated actor correctly returns to
\texttt{STARTUP\_RECOVERY}.

\subsection{Peripheral actors and discovery}

Keypads and sensors do not hard-code an \texttt{ActorRef} to the control unit.
Instead, they subscribe to the Pekko receptionist using the shared
\texttt{CONTROL\_UNIT\_SERVICE\_KEY}. This keeps the deployment \textbf{decoupled} from
specific node addresses.

The siren uses the same pattern in the opposite direction: it is registered as
a local service, and the control unit subscribes to siren updates so that it
can activate or deactivate the available siren actor.

\subsection{Recovery-oriented state machine}

The cluster version preserves the same \textbf{alarm semantics} as assignment 3:
\texttt{DISARMED}, \texttt{EXIT\_DELAY}, \texttt{ARMED},
\texttt{ENTRY\_DELAY}, and \texttt{ALARM}. The only important extension is the
safe startup state \texttt{STARTUP\_RECOVERY}.

After creation or recreation, the control unit does not guess whether the
system was previously armed or disarmed. It ignores or logs sensor events
until a correct PIN is received. This \textbf{safe recovery behavior} is directly
aligned with the failure-and-recovery requirement of the assignment.
